\documentclass[11pt,letterpaper]{article}
\usepackage[left=1.5in,right=1in,top=1in,bottom=1in]{geometry}
\usepackage{fancyhdr}
\usepackage{array}
\usepackage{longtable}
\usepackage{amssymb}
\usepackage{enumitem}

\pagestyle{fancy}
\fancyhf{}
\fancyhead[L]{\small Owner-Builder Job Site Binder | Section 2: Contracts \& Legal Documents | Page \thepage}
\fancyfoot[C]{\thepage}
\fancyfoot[R]{\small \copyright{} 2024 Build Your House}
\renewcommand{\headrulewidth}{0.4pt}
\renewcommand{\footrulewidth}{0.4pt}

\setlength{\parindent}{0pt}
\setlength{\parskip}{6pt}

\begin{document}

\begin{center}
\textbf{\Large DISPUTE RESOLUTION PROCEDURE}
\end{center}

\vspace{12pt}

\textbf{Project Name:} \rule{8cm}{0.4pt}

\textbf{Project Address:} \rule{8cm}{0.4pt}

\textbf{Owner:} \rule{9cm}{0.4pt}

\vspace{6pt}

\textit{This document outlines the procedures for resolving disputes between the owner and contractors/subcontractors. Following these steps can help resolve issues efficiently and cost-effectively while maintaining professional relationships.}

\vspace{12pt}

\section*{STEP 1: INITIAL COMMUNICATION}

\textbf{Timeline:} Within 3 business days of issue arising

\textbf{Objective:} Resolve the issue through direct, informal communication

\vspace{6pt}

\textbf{Actions:}

\begin{enumerate}
\item Identify the specific issue or concern
\item Contact the other party directly (phone call or in-person meeting preferred)
\item Discuss the issue calmly and professionally
\item Listen to the other party's perspective
\item Attempt to reach a mutually acceptable solution
\item Document the conversation and any agreements reached
\end{enumerate}

\vspace{6pt}

\textbf{Documentation:}

Date of communication: \rule{4cm}{0.4pt}

Method: $\square$ Phone $\square$ In-person $\square$ Email $\square$ Text

Parties involved: \rule{8cm}{0.4pt}

Summary of discussion:

\rule{14cm}{0.4pt}

\rule{14cm}{0.4pt}

\rule{14cm}{0.4pt}

Agreement reached: $\square$ Yes $\square$ No

If yes, describe: \rule{9cm}{0.4pt}

\rule{14cm}{0.4pt}

\vspace{6pt}

\textbf{If issue is resolved at this step, no further action is needed. If not resolved, proceed to Step 2.}

\newpage

\section*{STEP 2: WRITTEN NOTICE}

\textbf{Timeline:} Within 7 business days if Step 1 does not resolve the issue

\textbf{Objective:} Formally document the dispute and proposed resolution

\vspace{6pt}

\textbf{Actions:}

\begin{enumerate}
\item Prepare written notice of dispute including:
   \begin{itemize}
   \item Detailed description of the issue
   \item Relevant contract provisions or agreements
   \item Timeline of events
   \item Supporting documentation (photos, invoices, emails, etc.)
   \item Desired resolution
   \item Deadline for response (typically 7 business days)
   \end{itemize}
\item Deliver notice via certified mail, return receipt requested
\item Keep copies of all documentation
\item Await written response from other party
\end{enumerate}

\vspace{6pt}

\textbf{Written Notice Checklist:}

$\square$ Description of dispute clearly stated

$\square$ Relevant dates and timeline included

$\square$ Supporting documents attached

$\square$ Specific resolution requested

$\square$ Response deadline specified

$\square$ Sent via certified mail

$\square$ Copy kept for records

\vspace{6pt}

\textbf{Documentation:}

Date notice sent: \rule{4cm}{0.4pt}

Certified mail tracking \#: \rule{5cm}{0.4pt}

Date response received: \rule{4cm}{0.4pt}

Response summary:

\rule{14cm}{0.4pt}

\rule{14cm}{0.4pt}

\rule{14cm}{0.4pt}

Agreement reached: $\square$ Yes $\square$ No

\vspace{6pt}

\textbf{If issue is resolved at this step, document the agreement in writing and have both parties sign. If not resolved, proceed to Step 3.}

\newpage

\section*{STEP 3: FORMAL MEETING}

\textbf{Timeline:} Within 14 business days of written notice if issue remains unresolved

\textbf{Objective:} Meet face-to-face to negotiate a resolution

\vspace{6pt}

\textbf{Actions:}

\begin{enumerate}
\item Schedule meeting at mutually convenient time and neutral location
\item Both parties should bring:
   \begin{itemize}
   \item All relevant documentation
   \item Contract and any change orders
   \item Photos, invoices, correspondence
   \item Witnesses (if applicable)
   \end{itemize}
\item Conduct meeting professionally:
   \begin{itemize}
   \item Present facts and evidence
   \item Allow each party to speak without interruption
   \item Focus on solutions, not blame
   \item Explore compromise options
   \end{itemize}
\item Document meeting minutes and any agreements
\item Both parties sign meeting summary
\end{enumerate}

\vspace{6pt}

\textbf{Meeting Documentation:}

Date of meeting: \rule{4cm}{0.4pt} Time: \rule{3cm}{0.4pt}

Location: \rule{10cm}{0.4pt}

Attendees: \rule{11cm}{0.4pt}

\rule{14cm}{0.4pt}

Summary of discussion:

\rule{14cm}{0.4pt}

\rule{14cm}{0.4pt}

\rule{14cm}{0.4pt}

\rule{14cm}{0.4pt}

\rule{14cm}{0.4pt}

Proposed solutions discussed:

\rule{14cm}{0.4pt}

\rule{14cm}{0.4pt}

\rule{14cm}{0.4pt}

Agreement reached: $\square$ Yes $\square$ No

\vspace{6pt}

\textbf{If agreement is reached, both parties must sign a written settlement agreement. If not resolved, proceed to Step 4.}

\newpage

\section*{STEP 4: MEDIATION}

\textbf{Timeline:} Within 30 days if Steps 1-3 do not resolve the issue

\textbf{Objective:} Use neutral third-party mediator to facilitate resolution

\vspace{6pt}

\textbf{What is Mediation?}

Mediation is a voluntary process where a neutral third party (mediator) helps both sides reach a mutually acceptable agreement. The mediator does not make decisions but facilitates communication and negotiation.

\vspace{6pt}

\textbf{Benefits of Mediation:}

\begin{itemize}[leftmargin=0.5in]
\item Less expensive than arbitration or litigation
\item Faster resolution (typically 1-2 sessions)
\item Confidential process
\item Preserves business relationships
\item Both parties control the outcome
\item Non-binding (unless agreement is reached)
\end{itemize}

\vspace{6pt}

\textbf{Actions:}

\begin{enumerate}
\item Both parties agree to participate in mediation
\item Select mediator:
   \begin{itemize}
   \item Mutually agreed upon individual
   \item Local mediation service
   \item Construction dispute specialist
   \end{itemize}
\item Split mediation costs equally (unless otherwise agreed)
\item Prepare mediation statement outlining position and desired outcome
\item Attend mediation session(s)
\item If agreement is reached, sign binding settlement agreement
\end{enumerate}

\vspace{6pt}

\textbf{Mediation Resources:}

\begin{itemize}[leftmargin=0.5in]
\item American Arbitration Association (AAA): www.adr.org
\item Local bar association mediation services
\item State court-annexed mediation programs
\item Private mediation services
\end{itemize}

\vspace{6pt}

\textbf{Mediation Documentation:}

Mediator selected: \rule{7cm}{0.4pt}

Contact: \rule{7cm}{0.4pt} Phone: \rule{3cm}{0.4pt}

Date of mediation: \rule{4cm}{0.4pt}

Cost: \$ \rule{3cm}{0.4pt} (split equally)

Agreement reached: $\square$ Yes $\square$ No

Settlement amount/terms: \rule{7cm}{0.4pt}

\rule{14cm}{0.4pt}

\vspace{6pt}

\textbf{If mediation is successful, both parties sign settlement agreement and dispute is resolved. If mediation fails, proceed to Step 5.}

\newpage

\section*{STEP 5: ARBITRATION OR LITIGATION}

\textbf{Timeline:} If all other methods fail

\textbf{Objective:} Obtain binding decision from arbitrator or court

\vspace{6pt}

\textbf{OPTION A: BINDING ARBITRATION}

\textbf{What is Arbitration?}

Arbitration is a formal process where a neutral arbitrator (like a private judge) hears evidence from both sides and makes a binding decision. The decision is final and generally cannot be appealed.

\vspace{6pt}

\textbf{Process:}

\begin{enumerate}
\item File demand for arbitration with arbitration organization (AAA, JAMS, etc.)
\item Pay filing fee (typically \$1,000-\$3,000+)
\item Select arbitrator (or use arbitration service's selection process)
\item Conduct discovery (exchange of information/evidence)
\item Attend arbitration hearing
\item Arbitrator issues binding decision
\item Decision is enforceable in court
\end{enumerate}

\vspace{6pt}

\textbf{Pros:}
\begin{itemize}[leftmargin=0.75in]
\item Faster than court litigation
\item Less formal procedures
\item Arbitrator may have construction expertise
\item Private and confidential
\end{itemize}

\textbf{Cons:}
\begin{itemize}[leftmargin=0.75in]
\item Can be expensive (arbitrator fees, attorney fees)
\item Limited appeal rights
\item Discovery may be limited
\end{itemize}

\vspace{6pt}

\textbf{OPTION B: SMALL CLAIMS COURT}

\textbf{When to Use:}

Small claims court is appropriate when the dispute amount is below the state's small claims limit (typically \$5,000-\$10,000 depending on state).

\vspace{6pt}

\textbf{Process:}

\begin{enumerate}
\item File claim at county small claims court
\item Pay filing fee (typically \$30-\$100)
\item Serve defendant with claim
\item Attend hearing (no attorneys required in most states)
\item Present evidence to judge
\item Judge issues decision (usually immediately)
\end{enumerate}

\vspace{6pt}

\textbf{Pros:}
\begin{itemize}[leftmargin=0.75in]
\item Inexpensive and fast
\item Simple procedures, no attorney needed
\item Quick resolution (typically within 30-90 days)
\end{itemize}

\textbf{Cons:}
\begin{itemize}[leftmargin=0.75in]
\item Limited to lower dollar amounts
\item Limited discovery
\item Appeal rights may be limited
\end{itemize}

\vspace{6pt}

\textbf{OPTION C: CIVIL LITIGATION}

\textbf{When to Use:}

Civil litigation in regular court is typically used for larger disputes exceeding small claims limits or when arbitration is not required by contract.

\vspace{6pt}

\textbf{Process:}

\begin{enumerate}
\item Hire attorney
\item File complaint in civil court
\item Serve defendant with complaint
\item Defendant files answer
\item Discovery phase (depositions, interrogatories, document requests)
\item Pre-trial motions and hearings
\item Trial (before judge or jury)
\item Judge or jury renders verdict
\item Right to appeal
\end{enumerate}

\vspace{6pt}

\textbf{Pros:}
\begin{itemize}[leftmargin=0.75in]
\item Full discovery rights
\item Right to jury trial
\item Right to appeal
\item Established rules of evidence and procedure
\end{itemize}

\textbf{Cons:}
\begin{itemize}[leftmargin=0.75in]
\item Expensive (attorney fees, court costs, expert witnesses)
\item Time-consuming (can take 1-3+ years)
\item Public record
\item Uncertain outcome
\end{itemize}

\vspace{12pt}

\section*{DOCUMENTATION REQUIREMENTS}

\textbf{Throughout the dispute resolution process, maintain documentation of:}

$\square$ All communications (emails, letters, text messages)

$\square$ Meeting notes and summaries

$\square$ Photos of disputed work or conditions

$\square$ Invoices, payment records, receipts

$\square$ Contract and any change orders

$\square$ Relevant building codes or specifications

$\square$ Expert reports or inspections

$\square$ Timeline of events

$\square$ Witness statements

$\square$ Any other relevant evidence

\vspace{6pt}

\textbf{Documentation storage location:} \rule{8cm}{0.4pt}

\vspace{12pt}

\section*{TIMELINE SUMMARY}

\begin{tabular}{|p{1.5in}|p{1in}|p{3.5in}|}
\hline
\textbf{Step} & \textbf{Timeline} & \textbf{Action} \\
\hline
1. Initial Communication & 3 days & Direct communication to resolve issue \\
\hline
2. Written Notice & 7 days & Formal written notice of dispute \\
\hline
3. Formal Meeting & 14 days & Face-to-face negotiation meeting \\
\hline
4. Mediation & 30 days & Third-party mediation \\
\hline
5. Arbitration/Litigation & As needed & Binding arbitration or court \\
\hline
\end{tabular}

\vspace{12pt}

\textbf{IMPORTANT NOTES:}

\begin{itemize}[leftmargin=0.5in]
\item Attempt to resolve disputes at the earliest possible step
\item Keep all communications professional and documented
\item Consult with an attorney before proceeding to arbitration or litigation
\item Check your contract for specific dispute resolution requirements
\item Be aware of statute of limitations for filing claims in your state
\item Consider the cost of dispute resolution versus the amount in dispute
\end{itemize}

\end{document}
